\chapter*{Anexo A: Guía de reproducibilidad de la plataforma}
\addcontentsline{toc}{chapter}{Anexo A: Guía de reproducibilidad de la plataforma}
\label{cap:anexo_reproduccion}

\setcounter{section}{0}
\renewcommand{\thesection}{A.\arabic{section}}
\setcounter{code}{0}
\renewcommand{\thecode}{A.\arabic{code}}

Este anexo recoge los pasos con los que se reconstruyó la plataforma desde
cero en el Capítulo~\ref{cap:visualizacion}. No introduce ninguna decisión
de diseño nueva. Es una referencia operativa que complementa con la sintaxis
exacta de los comandos el despliegue descrito en el
Capítulo~\ref{cap:infraestructura}, y su versión completa y mantenida vive
en \path{docs/development/tfm_lakehouse_workflow.md}.

\section{Prerrequisitos y versiones fijadas}
\label{sec:anexo_reproduccion_prerrequisitos}

La plataforma necesita WSL2 con systemd activado, Docker, \texttt{kubectl},
\texttt{helm} y \texttt{uv} para el entorno Python del propio repositorio.
Todas las versiones de los servicios se fijan de forma explícita, de modo
que dos despliegues distintos produzcan el mismo sistema. La
Tabla~\ref{tab:anexo_versiones} las reúne.

\begin{table}[H]
  \centering
  \small
  \renewcommand{\arraystretch}{1.15}
  \begin{tabular}{
    >{\raggedright\arraybackslash}p{0.3\textwidth}
    >{\raggedright\arraybackslash}p{0.3\textwidth}
    >{\raggedright\arraybackslash}p{0.26\textwidth}}
    \hline
    \textbf{Componente} & \textbf{Versión fijada} & \textbf{Origen} \\
    \hline
    Kubernetes con k3s & v1.36.4+k3s1 & Instalador oficial de k3s \\
    MinIO & \textit{chart} 17.0.21 & Referencia OCI de Bitnami \\
    PostgreSQL & \textit{chart} 18.11.3 & Referencia OCI de Bitnami \\
    Airflow & \textit{chart} 1.22.0 & Repositorio de Apache Airflow \\
    Superset & \textit{chart} 0.22.8, aplicación 6.1.0 & Repositorio de Apache Superset \\
    Spark & 3.5.9 & Distribución oficial con Hadoop 3 \\
    Iceberg & \textit{runtime} 1.11.0 para Spark 3.5 & Maven Central \\
    Acceso a S3 & \texttt{hadoop-aws} 3.3.4 y SDK 1.12.262 & Maven Central \\
    \hline
  \end{tabular}
  \caption{Versiones fijadas de los componentes desplegados.}
  \label{tab:anexo_versiones}
\end{table}

Las rutas de código y de datos se montan en los pods desde el directorio de
trabajo del nodo. Por eso los dos ficheros de flujo contienen la ruta
absoluta del repositorio en la máquina de desarrollo, y esa ruta debe
editarse antes de entregarlos desde cualquier otra máquina. Es una carencia
de portabilidad conocida, recogida en el Anexo~\ref{cap:anexo_limitaciones}.

\section{Clúster y servicios base}
\label{sec:anexo_reproduccion_servicios}

El Código~\ref{cod:anexo_cluster} instala k3s con el modo de red
\texttt{host-gw}, necesario bajo WSL2, y crea el espacio de nombres
dedicado. Los servicios se instalan después con los valores versionados bajo
\path{infra/helm-values/}, tras crear sus credenciales como secretos de
Kubernetes que nunca se guardan en el repositorio.

\begin{code}[H]
\begin{lstlisting}[language=bash, basicstyle=\footnotesize\ttfamily, breaklines=true]
curl -sfL https://get.k3s.io | \
  INSTALL_K3S_EXEC="--flannel-backend=host-gw --write-kubeconfig-mode=644" sh -
kubectl create namespace tfm-lakehouse

helm install minio oci://registry-1.docker.io/bitnamicharts/minio \
  --version 17.0.21 -f infra/helm-values/minio-values.yaml -n tfm-lakehouse
helm install postgresql oci://registry-1.docker.io/bitnamicharts/postgresql \
  --version 18.11.3 -f infra/helm-values/postgresql-values.yaml -n tfm-lakehouse
helm install airflow apache-airflow/airflow \
  --version 1.22.0 -f infra/helm-values/airflow-values.yaml -n tfm-lakehouse
\end{lstlisting}
\caption{Instalación del clúster y de los tres servicios base.}
\label{cod:anexo_cluster}
\end{code}

\section{Imágenes de Spark y de ingesta}
\label{sec:anexo_reproduccion_imagenes}

El flujo de transformación ejecuta sus tres trabajos sobre una única
imagen, construida por capas a partir de la imagen oficial de Spark. El
Código~\ref{cod:anexo_imagenes} muestra la construcción y la importación de
esa imagen y de la imagen de ingesta. La importación en k3s es el único paso
que exige privilegios de administrador.

\begin{code}[H]
\begin{lstlisting}[language=bash, basicstyle=\footnotesize\ttfamily, breaklines=true]
docker build -f infra/spark-conf/Dockerfile.gold \
  -t tfm-lakehouse/spark-py:3.5.9-iceberg1.11.0-gold .
docker save tfm-lakehouse/spark-py:3.5.9-iceberg1.11.0-gold | \
  sudo k3s ctr images import -

docker build -f infra/ingest/Dockerfile -t tfm-lakehouse/ingest:py3.14 .
docker save tfm-lakehouse/ingest:py3.14 | sudo k3s ctr images import -

kubectl apply -f infra/spark-conf/spark-rbac.yaml
\end{lstlisting}
\caption{Construcción e importación de las imágenes de la plataforma.}
\label{cod:anexo_imagenes}
\end{code}

\section{Ejecución de los flujos}
\label{sec:anexo_reproduccion_flujos}

El subconjunto masivo de ORCID exige descargar una sola vez el archivo de
resúmenes de perfil, de unos 46 GB, y verificar su suma de comprobación. La
tarea de ingesta lo extrae después por su cuenta y nunca lo descarga. Los
dos flujos se entregan copiando su fichero al volumen de definiciones de
Airflow y se lanzan con la línea de órdenes de Airflow, como muestra el
Código~\ref{cod:anexo_flujos}.

\begin{code}[H]
\begin{lstlisting}[language=bash, basicstyle=\footnotesize\ttfamily, breaklines=true]
P=$(kubectl get pod -n tfm-lakehouse -l component=dag-processor -o name | head -1)
kubectl cp dags/ingest_validate.py \
  tfm-lakehouse/${P#pod/}:/opt/airflow/dags/ingest_validate.py -c dag-processor
kubectl exec -n tfm-lakehouse ${P#pod/} -c dag-processor -- \
  airflow dags trigger ingest_validate --run-id run-1 \
  --conf '{"count": 1000, "bulk_max_records": 20000}'

kubectl cp dags/transform_publish.py \
  tfm-lakehouse/${P#pod/}:/opt/airflow/dags/transform_publish.py -c dag-processor
kubectl exec -n tfm-lakehouse ${P#pod/} -c dag-processor -- \
  airflow dags trigger transform_publish --run-id run-1
\end{lstlisting}
\caption{Entrega y lanzamiento de los dos flujos de Airflow.}
\label{cod:anexo_flujos}
\end{code}

Una ejecución con los parámetros por defecto tarda entre seis y diez
minutos con el entorno ya caliente. Los parámetros del flujo de ingesta
controlan el número de currículos sintéticos, la semilla, la proporción de
currículos enlazados, el tamaño de la muestra de la API y el límite de
registros masivos aterrizados. Los del flujo de transformación controlan el
número de ejecutores, su memoria, las particiones de mezcla y los umbrales
de rechazo y de similitud entre organizaciones.

\section{Panel de indicadores y verificación}
\label{sec:anexo_reproduccion_panel}

El panel se instala con la imagen propia de Superset, un rol de solo lectura
sobre la base de datos de gold y la importación del panel exportado, todo
ello descrito en \path{infra/superset/README.md}. El acceso se hace con una
redirección de puerto local. La verificación del conjunto combina la batería
de pruebas del repositorio, que se lanza con el comando documentado del
Código~\ref{cod:anexo_verificacion}, y la comprobación del estado real de
los pods y de los dos flujos.

\begin{code}[H]
\begin{lstlisting}[language=bash, basicstyle=\footnotesize\ttfamily, breaklines=true]
uv sync --group codegen --group testing
uv run pytest -n auto tests

kubectl get pods -n tfm-lakehouse
kubectl port-forward -n tfm-lakehouse svc/superset 8088:8088
\end{lstlisting}
\caption{Verificación del repositorio y del estado del clúster.}
\label{cod:anexo_verificacion}
\end{code}

\section{Recuperación de los fallos conocidos}
\label{sec:anexo_reproduccion_recuperacion}

Dos fallos requieren una intervención manual documentada, ambos recogidos en
el registro del Anexo~\ref{cap:anexo_limitaciones}.

\begin{itemize}
  \item \textbf{Reinicio bloqueado del planificador de Airflow.} Una fila de
  tarea huérfana que dejó un planificador caído hace fallar todos los reinicios
  posteriores. Se resuelve marcando como fallidas esa tarea y su ejecución
  desde un pod sano, borrando el pod que nombra si aún existe y reiniciando el
  pod del planificador.
  \item \textbf{Usuario de PostgreSQL sin crear.} Si la publicación falla
  con un error de autenticación del usuario de gold tras un arranque en
  frío, el primer arranque de la base de datos quedó interrumpido. Se
  confirma revisando el registro del contenedor anterior y se resuelve
  borrando el pod y su volumen persistente para forzar una inicialización
  limpia.
\end{itemize}
